% Nouveau EGA V (cote 25): archive cote 25, site page 60 (PDF page 61).
% typescript. Draft compared: none.
% AI-transcribed by Claude Opus 5.5 (claude-opus-5-5, adaptive thinking, effort high), Message Batches API, batch ega25b (msgbatch_01A7n5YuXubUPNK1zsepi8RG), 2026-11-13; not yet checked by eye.
\scanpage{25/60}
 \mspage{55}

Pour comprendre la signification de ce théorème, il conviendrait

de noter dès le début du nº \hand{(}que  si S=Spec(k)\hand{)}, alors pour une droite
L ``assez générale'' dans $\hand{\check{\ }}P$ ,  la condition de régularité est vérifiée (\unsure{notamment pour S.3.|notamment pour 5.3.})
(cf ratrappages signalés au  nº\unsure{l'avant}).  Par suite, dans les construc%
tions de ``bons'' pinceaux linéaires  indiquées au début du nº, on pouvait
imposer que les pinceaux envisagés  satisfaisaient ladite conditi-
on  (ce qui  est une conditionde même type mais\hand{,} distincte de celle qui consiste à
exiger que pour tout \hand{$\zeta\in$}L,  $\phi_{\hand{\zeta}}$  soit $X_{k(\hand{\zeta})}$-régulier ). Il fallait inclu
re la condition en question dans les énoncés récapitulatifs préconisés.
  D'autre part, pratiquement  10.2. est
utilisé surtout dans la situation 9.5. ,  ce qui  rend  désirable de
n'énoncer la réformulation de  9.5.  en termes de pinceaux qu'après
10.2. , pour pouvoir  inclure dans l'énoncé en question  également l'iso
morphisme du pinceau avec un éclaté (i.e. donner une description de la
situation  permettant une référence commode).  On obtient ainsi un
moyen, pour tout schéma projectif lisse \hand{(}géom  conn  X  de dim$\geq$2 \hand{)}sur un corps infini k,  de trou
ver un sous-schéma fermé lisse non vide de codimension 2 en chacun de

ses points 
 tel que le schéma éclaté admette une fibration
sur P$^1$ , dont toutes les\hand{/}fibres \hand{(}sont géom intègres, et  dont toutes les fibres \hand{)}sont lisses sont au plus un nombre fini,
ces dernières ayant au plus un point  géom. singulier, et
un tel point  étant rationnel sur k et géom. singulier ordinaire .  Cela
explique l'importance  de l'étude approfondie ( tout juste abordée à l'heure
actuelle) de telles fibrations à fibres singulières, pour réduire dans 
une certaine mesure l'étude des variétés projectives lisses de dim d
à celle des(familles à un paramètre de) variétés projectives de dim d-1\hand{,}
\hand{pouvant avoir des singularités ordinaires.}

\placedmargin{60}
